\documentclass[11pt]{article}
\usepackage{fontspec}
\usepackage[french,english]{babel}
\usepackage{amsmath}
\usepackage{amssymb}
\usepackage{amsthm}
\usepackage{mathrsfs}
\usepackage{tikz-cd}
\usepackage{enumitem}

\newcommand{\printedpage}[1]{\marginpar{\footnotesize p.~#1}}
\newcommand{\msp}[1]{\marginpar{\footnotesize [#1]}}
\newcommand{\Dl}{\boldsymbol{\Delta}}
\newcommand{\tl}{\tilde{}}
\newcommand{\Ob}{\mathrm{Ob}}
\newcommand{\Fl}{\mathrm{Fl}}

\begin{document}

\printedpage{49}
% head: A. Grothendieck : Dérivateurs, Ch. XVII, Édition des Universités de Montpellier II et Paris VII
% foot: Version 9 février 2019

\msp{page 98}

\section*{6\quad Application to the construction of derivators (associated with special Quillen triples)}

\textbf{Theorem 6.1.}\footnote[29]{This is the main result in this part XVII of the notes.} \textit{Let $(W,C,F)$ be a special closed Quillen triple on the category $\mathcal{M}$} (cf. def. 5.6, p. 90, and cor. 5.7). \textit{For every $I$ in} Cat\textit{, let us consider the category}
\[
\mathcal{M}(I) \;=\; \underline{\mathrm{Hom}}(I^{\circ},\mathcal{M}) \;, \tag*{(6.1.1)}
\]
\textit{and let us consider}
\[
W(I),\ F(I) \;\subset\; \Fl(\mathcal{M}(I)) \;,
\]
\textit{defined by}
\[
\begin{aligned}
&\text{\textit{$u : X \longrightarrow Y$ in $\mathcal{M}(I)$ is in $W(I)$ (resp. $F(I)$) if and only}}\\
&\text{\textit{if $\forall\, i \in \Ob I$, $u(i) : X(i) \longrightarrow Y(i)$ is in $W$ (resp. $F$).}}
\end{aligned}
\tag*{(6.1.2)}
\]

\begin{itemize}[leftmargin=2.5em]
\item[a)] \textit{There exists a closed Quillen structure (necessarily unique) in $\mathcal{M}(I)$, having $W(I)$ and $F(I)$ as the sets of quasi-isomorphisms and of fibrations respectively. We thus also have}
\[
\left\{
\begin{aligned}
TF(I) \;&=\; F(I)\cap W(I) \;=\; \{u\in \Fl(\mathcal{M}(I)) \mid u(i)\in TF \ \ \forall\, i\in \Ob I\},\\
C(I) \;&=\; TF(I)^{*}.
\end{aligned}
\right.
\tag*{(6.1.3)}
\]
\textit{This Quillen structure is} special.

\msp{page 99}
\item[b)] \textit{Let $I_0$ be the discrete category having the same set of objects as $I$, and}
\[
v_I \ \ \text{\textit{or}}\ \ v \;:\; I_0 \longrightarrow I
\tag*{(6.1.4)}
\]
\textit{the canonical functor. Let us denote by}
\[
C(I) \;=\; (TF(I))^{*}, \qquad TC(I) \;=\; C(I)\cap W(I) \;=\; F(I)^{*}
\tag*{(6.1.5)}
\]
\textit{the sets of cofibrations resp. of trivial cofibrations in $\mathcal{M}(I)$, and let us use the same notations $C(I_0)$, $TC(I_0)$ for $I_0$. We then have}
\[
\begin{aligned}
&\text{\textit{$u : X_0 \longrightarrow Y_0$ in $\mathcal{M}(I_0)$ is in $C_0$ (resp. $TC_0$) if and only if}}\\
&\text{\textit{for every $i$ in $\Ob I$, $u(i) : X_0(i) \longrightarrow Y_0(i)$ is in $C$ (resp. $TC$).}}
\end{aligned}
\tag*{(6.1.6)}
\]
\textit{Finally we have}
\[
\left\{
\begin{aligned}
C(I) \;&=\; (v_!(C(I_0)))\tl\\
TC(I) \;&=\; (v_!(TC(I_0)))\tl\,.
\end{aligned}
\right.
\tag*{(6.1.7)} % ?? uncertain: shape of the small superscript mark after the closing parenthesis -- tilde (\tilde{}) or breve (\breve{})
\]

\item[c)] \textit{Let}
\[
u \;:\; I \longrightarrow I'
\]
\textit{be an arrow in} Cat\textit{, let us consider the couple of adjoint functors}
\[
\begin{tikzcd}
\mathcal{M}(I) \arrow[r, shift left, "u_!"] & \mathcal{M}(I') \arrow[l, shift left, "u^{*}"]
\end{tikzcd}
\;.
\tag*{(6.1.8)}
\]

\printedpage{50}
% head: A. Grothendieck : Dérivateurs, Ch. XVII, Édition des Universités de Montpellier II et Paris VII
% foot: Version 9 février 2019

\msp{page 100}
\textit{This couple satisfies Quillen's conditions for couples of adjoint functors between Quillen categories} (Hom. Alg. Chp. I, §4, th. 3, p. 4.5), \textit{so that they define by passage to the localized categories a couple of adjoint functors}
\[
\begin{tikzcd}
\mathrm{Ho}_W(I) \arrow[r, shift left, "u_!^{W}"] & \mathrm{Ho}_W(I') \arrow[l, shift left, "u^{*}_{W}"]
\end{tikzcd}
\;.
\tag*{(6.1.9)}
\]
\textit{We have, more precisely, the relations}
\[
\left\{
\begin{aligned}
u_!(C(I)) \;&\subset\; C(I')\\
u_!(TC(I)) \;&\subset\; TC(I')\;,
\end{aligned}
\right.
\tag*{(6.1.10)}
\]
\textit{whence}
\[
(u_!(\mathcal{M}(I)_{\mathrm{c}})) \;\subset\; \mathcal{M}(I')_{\mathrm{c}}\;,
\tag*{(6.1.11)}
\]
\textit{and moreover}
\[
\begin{aligned}
&\text{\textit{for every arrow $f$ in $\mathcal{M}(I)_{\mathrm{c}}$, such that $f \in W(I)$, we have}}\\
&\text{\textit{$u_!(f) \in W(I')$, i.e. $(u_!)_{\mathrm{c}} : \mathcal{M}(I)_{\mathrm{c}} \longrightarrow \mathcal{M}(I')_{\mathrm{c}}$ is compatible}}\\
&\text{\textit{with the localizers.}}
\end{aligned}
\tag*{(6.1.12)}
\]
\textit{On the other hand we have}
\[
\left\{
\begin{aligned}
u^{*}(F(I')) \;&\subset\; F(I)\\
u^{*}(TF(I')) \;&\subset\; TF(I)\\
u^{*}(W(I')) \;&\subset\; W(I)\;,
\end{aligned}
\right.
\tag*{(6.1.13)}
\]
\textit{hence}
\[
\left\{
\begin{aligned}
&u^{*}(\mathcal{M}(I')_{\mathrm{f}}) \;\subset\; \mathcal{M}(I)_{\mathrm{f}}\\
&u^{*}_{\mathrm{f}} : \mathcal{M}(I')_{\mathrm{f}} \longrightarrow \mathcal{M}(I)_{\mathrm{f}} \text{ \textit{is compatible with the localizers.}}
\end{aligned}
\right.
\tag*{(6.1.14)}
\]

\msp{page 101}
\item[d)] \textit{Suppose that $I$ is an ordered category, and let us consider the set $C_I \subset \Fl(\mathcal{M}(I))$ defined in §2}\footnote[30]{In paragraph 2, this was denoted $C_{I^{\circ}}$ instead of $C_I$.} (cf. page 25, cf. also th. 2.1, p. 26, and its corollaries). \textit{We then have}
\[
\left\{
\begin{aligned}
C_I \;&\supset\; C(I)\\
TC_I \;&\supset\; TC(I)\;,
\end{aligned}
\right.
\tag*{(6.1.15)}
\]
\textit{and we have equality when $I$ is an ordered set that is} Artinian.
\end{itemize}

\textsc{Proof}. Assertion a) is a special case of 5.8, applied to the couple of adjoint functors
\[
\begin{tikzcd}
\mathcal{M}(I_0) \arrow[r, shift left, "v_!"] & \mathcal{M}(I) \arrow[l, shift left, "v^{*}"]
\end{tikzcd}
\;.
\]
The conditions (5.8.2) are satisfied, this was seen in (4.8). Proposition 5.8 then tells us of the existence of a closed Quillen triple structure on $\mathcal{M}(I)$, whose set of quasi-isomorphisms $W(I)$ and set of fibrations $F(I)$, are given by the first formula (5.8.3) and by the first formula (5.8.4), which is equivalent to the definitions (6.1.2). (\textbf{N.B.} The fact that $\mathcal{M}(I_0)$, with the obvious structure $W(I_0)$, $C(I_0)$, $F(I_0)$ (product), is
\printedpage{51}
% head: A. Grothendieck : Dérivateurs, Ch. XVII, Édition des Universités de Montpellier II et Paris VII
% foot: Version 9 février 2019
\msp{page 102}
indeed a special closed Quillen category, is $\pm$ tautological.)

\textit{Let us prove} b). The relation (6.1.6) is tautological. The relations (6.1.7) are contained in 5.8 (cf. second relation in (5.8.3) for $C' = v_!(C)\tl$, i.e. $C(I) = (v_!(C(I_0)))\tl$, and (5.8.5) for $TC(I) = (v_!(TC(I_0)))\tl$).

\textit{Let us prove} c). The first assertion is, by \textsc{Quillen}, a special case of the relations (6.1.10) to (6.1.14), it is these that have to be established. The relations (6.1.13) and consequently (6.1.14) are tautological by virtue of the definitions, so it is a matter of establishing \textit{(6.1.10)} and (6.1.12).

For the formulas (6.1.10), I use the commutative diagram
\[
\begin{tikzcd}[row sep=large, column sep=large]
I \arrow[d, "u"'] & I_0 \arrow[l, "v"'] \arrow[d, "u_0"] \\
I' & I'_0\;, \arrow[l, "v'"']
\end{tikzcd}
\]
and I note that the functors $u_!^{\mathcal{M}}$, $u_{0\,!}^{\mathcal{M}}$, commuting with small direct limits,
\msp{page 103}
behave well with respect to the operation
\[
\Phi \longmapsto \Phi\tl\,,
\]
when $\Phi$ is a subset of $\Fl(\mathcal{M}(I))$ resp. of $\Fl(\mathcal{M}(I_0))$ which is Q-accessible, i.e. such that there exists a \textit{small} subset $\Phi_0 \subset \Phi$ with $\Phi \subset \tilde{\Phi}_0$, more precisely, we have
\[
u_!(\Phi\tl) \;\subset\; (u_!(\Phi))\tl\,.
\]
This noted, we have by b)
\[
C(I) \;=\; (v_!(C(I_0)))\tl\,,
\]
whence
\[
u_!(C(I)) \;\subset\; (\underbrace{u_!(v_!(C(I_0)))}_{(uv)_!(C(I_0))})\tl\,,
\]
and the right-hand side can also be written
\[
(v'_!u_{0\,!}(C(I_0)))\tl \;\overset{?}{\subset}\; (v'_!(I'_0))\tl \;=\; C(I')\;. % [sic] (v'_!(I'_0)) -> (v'_!(C(I'_0))), the C( ) around I'_0 apparently omitted
\]
Thus everything comes down to showing that
\[
u_{0\,!}(C(I_0)) \;\subset\; C(I'_0)\;,
\tag*{(6.1.16)}
\]
which follows at once from the computation of $u_{0\,!}$ by means of direct sums, and from the fact that $C$ is stable under arbitrary direct sums. This thus proves $u_!(C(I)) \subset C(I')$, and we prove in the same way $u_!(TC(I)) \subset TC(I')$, which reduces in the same way to the special case
\[
u_{0\,!}(TC(I_0)) \;\subset\; TC(I'_0)\;.
\tag*{(6.1.17)}
\]

\printedpage{52}
% head: A. Grothendieck : Dérivateurs, Ch. XVII, Édition des Universités de Montpellier II et Paris VII
% foot: Version 9 février 2019

\msp{page 104}

It remains to prove (6.1.12), and we proceed as for 3.3 (p. 54), using the fact that an arrow $f$ in $\mathcal{M}(I)_{\mathrm{c}}$ which is in $W(I)$, can be written in $\mathcal{M}(I)_{\mathrm{c}}$ in the form $f = pi$, where $i \in TC(I)$, and where $p$ is a left inverse (retraction) of an arrow $j \in TC(I)$. We will therefore have $u_!(f) = u_!(p)u_!(i)$, and $u_!(i) \in TC(I')$ by (6.1.10), and $u_!(p)$ is a left inverse of $u_!(j) \in TC(I')$. We then have $u_!(i), u_!(j), u_!(p) \in W(I')$, whence $u_!(f) \in W(I')$, q.e.d.

Let us finally prove d). In view of the definition of $C_I$, to establish $C(I) \subset C_I$, it suffices to see this:

\textbf{Corollary 6.2.}\footnote[31]{\textbf{N.B.} This result generalizes, under the general conditions in 6.2.5.} \textit{Let $I$ be arbitrary in} Cat\textit{, let $U \subset V$ be two closed sets in $I$ (i.e. open sets in $I^{\circ}$), and let $f : X_{\bullet} \longrightarrow Y_{\bullet}$ be an arrow in $C(I)$ (resp. in $TC(I)$). Let us also put}
\msp{page 105}
\[
X_U \;=\; \varinjlim_{i\in U} X(i), \qquad X_V \;=\; \varinjlim_{i\in U} X(i)\;, % [sic] second index i\in U -> i\in V, X_V should be the limit over V
\]
\textit{and likewise for $Y$, and let us consider the diagram}
\[
\begin{tikzcd}[row sep=large, column sep=large]
X_U \arrow[r, "\text{trans.}"] \arrow[d, "u_U"'] \arrow[dr, phantom, "\text{cocart.}"] & X_V \arrow[d] \arrow[dd, dotted, bend left=60, "u_V"] \\
Y_U \arrow[r] \arrow[dr, dotted, "\text{trans.}"'] & Y_{U,V} \arrow[d, "u_{U,V}"'] \\
 & Y_V\;.
\end{tikzcd}
\tag*{(6.2.1)}
\]
\textit{Then the arrow $u_{U,V}$ is in $C$ (resp. $TC$).}

\textsc{Proof} of 6.2. Let us admit that the restriction of an arrow of $C(I)$ resp. $TC(I)$ to a closed set $V$ of $I$ is in $C(V)$ resp. $TC(V)$ (cf. \dots{} below). This allows us to assume $V = I$. Giving $U$, which is an open set of $I^{\circ}$, is equivalent to giving an arrow $v : I^{\circ} \longrightarrow \Dl^1$ (taking the fiber at 0), or again to giving an arrow $u : I \longrightarrow (\Dl^1)^{\circ} \simeq \Dl^1$ (taking the fiber at 1). Now % [sic] le fibre -> la fibre, gender (cf. "la fibre en 0" just before)
\[
u \;:\; I \longrightarrow \Dl^1\;, \qquad U \;=\; I_1 \overset{\text{def}}{=} u^{-1}(\{1\})\;.
\tag*{(6.2.2)}
\]

\msp{page 106}
We then have
\[
X_U \;=\; u_!(X)(1), \quad X_I \;=\; u_!(X)(0)\;,
\]
and the transition morphism $X_U \longrightarrow X_I$ is deduced from the arrow $0 \xrightarrow{\ \partial\ } 1$ in $\Dl^1$, and likewise for $Y$. Let
\[
\xi \;=\; u_!(X), \qquad \eta \;=\; u_!(Y)\;,
\]
\printedpage{53}
% head: A. Grothendieck : Dérivateurs, Ch. XVII, Édition des Universités de Montpellier II et Paris VII
% foot: Version 9 février 2019
so the diagram (6.2.1) is identified with the diagram deduced from
\[
\varphi \;:\; \xi \longrightarrow \eta\;, \qquad \varphi \;=\; u_!(f)\;,
\]
as
\[
\begin{tikzcd}[row sep=large, column sep=large]
\xi(1) \arrow[r, "\xi(\partial)"] \arrow[d, "\varphi(1)"'] \arrow[dr, phantom, "\text{cocart.}"] & \xi(0) \arrow[d] \arrow[dd, dotted, bend left=60, "\varphi(0)"] \\
\eta(1) \arrow[r] \arrow[dr, dotted, "\eta(0)"'] & \tilde{\eta} \arrow[d, "\tilde{\varphi}"'] \\
 & \eta(0) \qquad .
\end{tikzcd}
\tag*{(6.2.3)}
\]
Now by 6.1.10 applied to $u : I \longrightarrow \Dl^1$, the relation $f \in C(I)$ resp. $f \in TC(I)$ implies $\varphi \in C(\Dl^1)$ resp. $\varphi \in TC(\Dl^1)$. Thus to establish 6.1, we are reduced to the special case where $I = \Dl^1$, $U = \{1\}$, $V = I$, hence to proving the

\textbf{Lemma 6.2.4.} \textit{Let $\varphi : \xi \longrightarrow \eta$ be an arrow in $\mathcal{M}(\Dl^1)$. If $\varphi \in C(\Dl^1)$ (resp.}
\msp{page 107}
\textit{$\varphi \in TC(\Dl^1)$), then the arrow $\tilde{\varphi}$ in the diagram} (6.2.3) \textit{is in $C$ (resp. $TC$).}

More generally, and more precisely:

\textbf{Lemma 6.2.5.} \textit{Let $\mathcal{M}$ be a category stable under amalgamated sums, and $\Phi$ a subset Q-closed on the left} (we will take $\Phi = C$, and $\Phi = TC$, in the case of lemma 6.2.4). \textit{Let $\Psi = \Phi_{*}$, and let us consider, for every $I$ in} Cat\textit{, $\Phi(I)$, $\Psi\langle I\rangle$ in $\mathcal{M}(I)$, defined by:}
\[
\begin{aligned}
&\text{\textit{$u : X \longrightarrow Y$ in $\mathcal{M}(I)$ is in $\Psi\langle I\rangle$ if and only if $\forall\, i \in \Ob I$,}}\\
&\text{\textit{$u(i) : X(i) \longrightarrow Y(i)$ is in $\Psi$,}}
\end{aligned}
\tag*{(6.2.5.1)}
\]
\[
\Phi(I) \;=\; \Psi\langle I\rangle^{*}.
\tag*{(6.2.5.2)}
\]

Let us take the case $I = \Dl^1 = \{0 \xrightarrow{\ \partial\ } 1\}$, and let $u : X \longrightarrow Y$ be an arrow in $\mathcal{M}(\Dl^1)$, corresponding to a commutative square
\[
\begin{tikzcd}[row sep=large, column sep=large]
X(1) \arrow[r, "X(\partial)"] \arrow[d, "u(1)"'] & X(0) \arrow[d, "u(0)"] \\
Y(1) \arrow[r, "Y(\partial)"] & Y(0)\;.
\end{tikzcd}
\tag*{(6.2.5.3)}
\]

Let us consider the corresponding diagram

\end{document}
